\documentclass[11pt]{article}

\usepackage[margin=1in]{geometry}
\usepackage{booktabs}
\usepackage{amsmath}

\title{A computational model of the evolution of 
       honeybee waggle-dance communication}
\author{Grzegorz Chrupa{\l}a}
\date{\today}

\begin{document}

\maketitle

\section{Introduction}

Honeybee waggle dances provide spatial information about resources away from
the nest. One ecological hypothesis is that such communication is favored when
food is difficult to discover independently but remains useful after a worker
has found it. Empirical and theoretical studies suggest that the value of dance
information and recruitment depends on the spatiotemporal distribution of food
resources, including patch density, patch size, reward variability, and habitat
\cite{sherman_visscher_2002,dornhaus_chittka_2004,dornhaus_etal_2006,beekman_lew_2008,donaldson_matasci_dornhaus_2012,schurch_gruter_2014,price_gruter_2015}.
We begin with a deliberately simple model of directional communication on an
exposed horizontal comb, where dance angle can map directly onto food
direction.

\section{Method}

Colonies are the reproducing entities, while workers are behavioral agents.
Each colony has seven heritable mean traits: directional precision investment,
called directional bias in the simulation code and tables, which controls how
strongly a dancing worker's angle is centered on the intended direction;
receiver attention, which controls how likely workers are to follow an
available dance; sender transposition, which controls how much a dancer maps
food direction onto a gravity-referenced dance angle; receiver transposition,
which controls how much dance followers interpret the dance as
gravity-referenced rather than direct pointing; comb tilt and comb orientation,
which define the dance plane in the world; and search limit, which controls how
far a worker searches along its chosen direction before giving up.
Individual workers are sampled around the colony means using stable
within-colony variation. Each dance also includes transient production and
interpretation noise.

The directional-precision trait should be read as a costly investment in
producing directionally informative dances, not as a pure measure of neural
precision, motivation, or dance duration. In this model, $b$ sets the angular
concentration of a produced dance and contributes to the per-dance cost paid by
the colony. Dance cost is represented as a baseline cost plus a
precision-dependent cost. The baseline component captures time, metabolic
effort, or opportunity cost of dancing regardless of directional precision,
whereas the precision-dependent component captures the additional cost of
making the dance directionally informative. The reported simulations set the
baseline dance cost to 0 and therefore include only the precision-dependent
component. Thus the results show the evolution of costly signal precision under
that simplified cost schedule, not the separate evolution of intrinsic accuracy
and signaling effort.

More concretely, a colony is parameterized by mean directional precision
investment $b \in [0,1]$, mean receiver attention $a \in [0,1]$, mean sender
transposition $s \in [0,1]$, mean receiver transposition $r \in [0,1]$, and
mean search limit $\ell \geq 0$. It also has a comb tilt $t \in [0,1]$ and a
comb orientation $\phi$. Orientation can be treated either as a circular
direction or, in constrained probes, as an axial variable where $\phi$ and
$\phi + 180^\circ$ represent the same comb plane. Each worker receives stable individual
deviations from these means, sampled from Gaussian noise with standard
deviation 0.08 for unit-interval traits. For search limit, the same scale is
multiplied by the maximum search distance. All traits are clipped to their
allowed ranges. These deviations represent within-colony worker variation, but
are not themselves heritable in the current model. At reproduction, only the
colony-level mean traits are inherited, with Gaussian mutation of standard
deviation 0.07, again scaled by maximum search distance for the search-limit
trait. Sender and receiver transposition mutations can be correlated by a
configurable coupling parameter; comb-orientation mutations use a separate
circular mutation scale. Comb tilt can also use its own mutation scale when an
experiment supplies a geometric precondition rather than asking tilt to evolve.

Each foraging episode generates a set of food sites in the world. A food site
has a direction, distance from the nest, angular width, value, and capacity.
Workers act sequentially within an episode. If at least one dance is available,
a worker follows a randomly chosen dance with probability given by its receiver
attention. Otherwise it searches independently in a uniformly random direction.
In both cases, the worker searches along the chosen direction up to its
individual search limit. A food site is found only if it lies within the
searched angular window and its distance is no greater than the worker's search
limit. Independent searchers that find food add a noisy dance for the discovered
site, so scouts are not assigned in advance: a scout is simply a worker that
found food independently. This makes both background discovery and recruitment
emergent properties of the food distribution and the evolving worker traits.
This design follows the logic of previous ecological models in which
recruitment benefits depend on the availability and discoverability of resource
patches rather than on communication having a fixed payoff advantage
\cite{dornhaus_etal_2006,beekman_lew_2008}.

Comb tilt and orientation define a three-dimensional comb plane. Direct
pointing projects the horizontal food direction onto that plane, so its
available cue strength depends on both the tilt and the orientation of the comb
relative to the food direction. Gravity-referenced mapping uses the projection
of gravity onto the comb plane together with an episode-level sun azimuth.
Gravity provides no in-plane reference on a horizontal comb and becomes a
stronger cue as the comb becomes vertical. The sender transposition trait blends
between direct pointing and sun/gravity-referenced encoding, and the receiver
transposition trait blends between the corresponding decoding rules. Each
episode samples the sun azimuth from a configurable daytime arc.

The dancer's angle is sampled from a circular von Mises distribution
centered on the encoded direction. The concentration parameter is $14b$, so
workers with low directional precision investment produce nearly random
angles, while workers with high directional precision investment produce
angles concentrated around the intended direction. Additional dance-production
noise with standard deviation 0.18 is added to the angle. A worker follows an
available dance with probability $a$; if it follows, it decodes the signaled
direction with interpretation noise of standard deviation 0.12. If it does not
follow a dance, it searches in a uniformly random direction. A worker succeeds
when its search direction falls within the angular width of an available food
site whose distance is no greater than the worker's search limit.

After many episodes, each colony receives a payoff from collected food minus
travel, communication, and attention costs. Daughter colonies are sampled in
proportion to payoff, and colony-level traits mutate by Gaussian noise. We
track the mean directional precision investment, receiver attention, sender
and receiver transposition, comb tilt, comb-orientation alignment, and search
limit, plus forager success rate across generations.

In each episode, collected food is the sum of the values of food sites found by
workers. The current simulations use food value 1.0. Workers also pay a travel
cost proportional to distance: successful workers pay the cost of reaching the
found food site, while unsuccessful workers pay the cost of searching up to
their limit. The default travel cost is 0.02 per distance unit, so even the
longest default trip remains profitable if it finds food, but long unsuccessful
searches are costly. The colony then pays a communication cost of
$c_0 + 0.02b_d$ for each dance added by an independent finder, where $c_0$ is
the baseline dance cost and $b_d$ is that dancer's directional precision
investment. In the experiments reported here, $c_0 = 0$. The colony also pays
an attention cost of 0.01 for each worker that follows a dance. Thus more
precise directional signaling, more dance-following, and longer search effort
can increase food collection, but they are not free. The current model
can also include a configurable benefit proportional to comb verticality, which
represents non-communication selective pressure for vertical combs. The
parameter
``workers per colony'' controls the pool from which active foragers are sampled,
and therefore the amount of within-colony worker heterogeneity, but payoff
opportunity is set more directly by the number of foraging episodes per colony
and the number of foraging attempts per episode.

\section{Results}

\subsection{Horizontal comb}

We compared three simple food distributions over ten random seeds each
(seeds 101--110). The
``hard'' condition had one narrow food site per episode, the baseline condition
had two moderately wide food sites, and the ``easy'' condition had many broad
food sites. Table~\ref{tab:food-distribution} reports the generation at which
mean directional precision investment first reached 0.30, plus final generation
trait and success values.

Each run used 60 colonies, 80 workers per colony, 30 generations, 50 foraging
episodes per colony per generation, 12 foraging attempts per episode, and
horizontal-comb tilt $t=0$. Food-site distances were sampled uniformly from
1.0 to 8.0, maximum search distance was 8.0, travel cost was 0.02 per distance
unit, baseline dance cost was 0, the precision-dependent dance-cost coefficient
was 0.02, and food-site value was 1.0. All three conditions used food-site
capacity 6. The conditions varied only the number and angular width of food
sites: hard used one site of width 0.08, baseline used two sites of width 0.20,
and easy used eight sites of width 0.50. The hard condition therefore makes
independent discovery unlikely, while the easy condition makes independent
discovery likely even without using the dance.

% Generated by experiments/run_report_artifacts.py; do not edit by hand.
% Result command: python -u experiments/run_report_artifacts.py results --only report_conditions
% Result file: results/report_conditions.csv
\begin{table}[h]
\centering
\small
\begin{tabular}{lrrrrrrr}
\toprule
Condition & Reached & Reach gen. & Prec. & Attention & Search & Success & Payoff \\
\midrule
Hard & 4/10 & 22.750 & 0.264 & 0.313 & 2.789 & 0.008 & 0.001 \\
Baseline & 10/10 & 7.700 & 0.833 & 0.869 & 7.215 & 0.192 & 0.690 \\
Easy & 5/10 & 19.600 & 0.279 & 0.369 & 7.067 & 0.708 & 7.366 \\
\bottomrule
\end{tabular}
\caption{Food-distribution experiment under the current dynamic recruitment
model. ``Reached'' is the number of seeds in which mean directional precision
investment reached 0.30 by generation 30. ``Reach gen.'' is averaged over
reached seeds only. Prec. is final mean directional precision investment;
attention, search limit, success, and payoff values are also final-generation
means over all ten seeds.}
\label{tab:food-distribution}
\end{table}


The result matches the ecological prediction qualitatively, but with a sharper
distinction than in the earlier fixed-scout model. In the hard environment,
food is so rare and narrow that independent discoveries seldom seed dances;
final success and payoff remain close to zero, and communication evolves only
in some seeds. In the baseline environment, independent discovery is sufficient
to seed recruitment cascades, and directional precision investment, receiver
attention, and search limit all rise strongly. In the easy environment, success
and payoff are high, but directional precision investment and receiver
attention remain much lower than in the baseline case because food is often
found without precise communication.

\subsection{Vertical comb}

We also compared colonies initialized with horizontal, tilted, and vertical
combs while holding the baseline food distribution fixed. Initial comb tilt was
set to $t=0$, $t=0.5$, or $t=1.0$, with ten seeds per condition. Comb tilt was
still heritable and mutable in these runs. Table~\ref{tab:comb-tilt} summarizes
the result after 30 generations.

% Generated by experiments/run_report_artifacts.py; do not edit by hand.
% Result command: python -u experiments/run_report_artifacts.py results --only report_conditions
% Result file: results/report_conditions.csv
\begin{table}[h]
\centering
\small
\begin{tabular}{lrrrrrrrr}
\toprule
Comb & Reached & Reach gen. & Prec. & Attention & Sender & Receiver & Success & Payoff \\
\midrule
Horizontal & 10/10 & 7.700 & 0.844 & 0.868 & 0.187 & 0.176 & 0.194 & 0.733 \\
Tilted & 10/10 & 8.100 & 0.829 & 0.886 & 0.171 & 0.196 & 0.195 & 0.727 \\
Vertical & 10/10 & 16.500 & 0.606 & 0.794 & 0.384 & 0.362 & 0.165 & 0.414 \\
\bottomrule
\end{tabular}
\caption{Comb-tilt comparison using the baseline food distribution and ten
seeds. Sender and receiver denote final mean sender and receiver transposition
traits.}
\label{tab:comb-tilt}
\end{table}


With mutable tilt, all three initial-tilt conditions reach the
directional-precision threshold in all ten seeds under the current baseline
settings. Tilted and vertical starts retain substantial foraging success over
30 generations, but they also evolve higher sender and receiver transposition
than horizontal starts. The initially vertical condition still has lower payoff
than the horizontal condition, and its final mean tilt is reduced in the
underlying result file, indicating that the short-run comparison does not by
itself establish a stable vertical-comb transition.

\subsection{Tilt-geometry calibration}

After replacing the scalar comb-tilt degradation with explicit comb-plane
geometry, we ran a calibration grid over initial comb tilt and the vertical-comb
benefit. These runs used the baseline food distribution, the daytime sun arc
from $0^\circ$ to $180^\circ$ centered at $90^\circ$, and ten seeds
(101--110). Success is the fraction of foraging attempts that found food.
Payoff is the final mean colony payoff after food rewards, search/travel costs,
dance costs, attention costs, and the vertical-comb benefit. Table~
\ref{tab:tilt-benefit-calibration} shows that increasing the vertical benefit
can preserve more vertical combs, but current settings produce a sharp tradeoff:
strong verticality often comes with reduced foraging success.

% Generated by experiments/run_report_artifacts.py; do not edit by hand.
% Result command: python -u experiments/run_report_artifacts.py results --only tilt_geometry_calibration
% Result file: results/tilt_geometry_calibration.csv
\begin{table}[h]
\centering
\small
\begin{tabular}{rrrrrrrr}
\toprule
Initial tilt & $\alpha$ & Final tilt & Sender & Receiver & Success & Payoff & Reached \\
\midrule
0.000 & 0.150 & 0.129 & 0.186 & 0.177 & 0.197 & 0.757 & 10/10 \\
0.000 & 0.200 & 0.136 & 0.188 & 0.205 & 0.193 & 0.731 & 10/10 \\
0.000 & 0.250 & 0.149 & 0.188 & 0.196 & 0.196 & 0.753 & 10/10 \\
0.500 & 0.150 & 0.154 & 0.183 & 0.196 & 0.192 & 0.708 & 10/10 \\
0.500 & 0.200 & 0.126 & 0.232 & 0.204 & 0.190 & 0.692 & 10/10 \\
0.500 & 0.250 & 0.237 & 0.211 & 0.250 & 0.187 & 0.672 & 10/10 \\
1.000 & 0.150 & 0.568 & 0.403 & 0.424 & 0.164 & 0.450 & 9/10 \\
1.000 & 0.200 & 0.292 & 0.248 & 0.261 & 0.176 & 0.580 & 10/10 \\
1.000 & 0.250 & 0.636 & 0.440 & 0.432 & 0.160 & 0.449 & 9/10 \\
\bottomrule
\end{tabular}
\caption{Ten-seed calibration of the proportional vertical-comb advantage under the explicit
tilt-geometry model. Sender and receiver are final mean transposition traits.
Reached is the number of seeds in which mean directional precision investment
reached 0.30 by generation 30.}
\label{tab:tilt-benefit-calibration}
\end{table}


The calibration suggests that the transition is path-dependent. Starting from a
horizontal comb, benefits in the range 0.15--0.25 do not move colonies far from
horizontal within 30 generations. Starting from a tilted comb, the same benefit
range can retain intermediate tilt but with lower success. Starting from a
vertical comb, benefits around 0.20--0.25 retain high verticality, but foraging
nearly fails in most seeds. A separate five-seed stress test showed that a much
larger benefit of 0.50 can drive high verticality from tilted starts, but at the
cost of severe loss of foraging success. Thus the current model can impose a
substantial architecture pressure, but the 30-generation calibration may end
before sender transposition, receiver transposition, and directional precision
have enough time to adapt together. The relevant open question is therefore
whether these cases represent a stable collapse or a longer coevolutionary
transition.

\subsection{Tilt orientation}

We also examined the evolved orientation of the comb tilt. The current sun arc
is centered at $90^\circ$, so Table~\ref{tab:orientation-calibration} reports
the weighted mean orientation in degrees and its offset from this sun-arc
center. The within-alignment column measures how strongly colonies within a run
converge on an orientation. The across-alignment column measures whether
different seeds converge on the same absolute orientation, weighting each seed
by its within-run orientation alignment.

% Generated by experiments/run_report_artifacts.py; do not edit by hand.
% Result command: python -u experiments/run_report_artifacts.py results --only orientation_calibration
% Result file: results/orientation_calibration.csv
\begin{table}[h]
\centering
\small
\begin{tabular}{rrrrrrrr}
\toprule
Initial tilt & Benefit & Final tilt & Success & Within align. & Mean orient. & Sun offset & Across align. \\
\midrule
0.000 & 0.150 & 0.142 & 0.195 & 0.518 & 122.841 & 32.8 & 0.473 \\
0.000 & 0.200 & 0.162 & 0.195 & 0.489 & 155.097 & 65.1 & 0.492 \\
0.000 & 0.250 & 0.151 & 0.194 & 0.478 & 79.011 & -11.0 & 0.084 \\
0.500 & 0.150 & 0.373 & 0.137 & 0.490 & 158.945 & 68.9 & 0.249 \\
0.500 & 0.200 & 0.445 & 0.137 & 0.356 & 214.421 & 124.4 & 0.158 \\
0.500 & 0.250 & 0.330 & 0.154 & 0.523 & 73.900 & -16.1 & 0.401 \\
1.000 & 0.150 & 0.737 & 0.090 & 0.545 & 251.668 & 161.7 & 0.374 \\
1.000 & 0.200 & 0.919 & 0.001 & 0.374 & 263.771 & 173.8 & 0.093 \\
1.000 & 0.250 & 0.917 & 0.002 & 0.392 & 266.776 & 176.8 & 0.291 \\
\bottomrule
\end{tabular}
\caption{Orientation sanity check under the explicit tilt-geometry model.
Mean orientation and sun offset are measured in degrees. Across alignment is
the weighted across-seed orientation alignment; values near 0 indicate that
different seeds choose different absolute orientations.}
\label{tab:orientation-calibration}
\end{table}


Orientation sometimes aligns within a run, but it does not yet show a clean,
robust relationship to the daytime sun arc across seeds. Some conditions point
near the sun-arc center, some near the opposite direction, and several have weak
across-seed alignment. This may indicate weak or path-dependent selection on
orientation. It also highlights a modeling issue: because a comb plane is
unchanged by reversing its normal, comb orientation may be better treated as an
axial variable, where $\phi$ and $\phi + 180^\circ$ are equivalent, rather than
as a fully circular direction.

\subsection{Constrained vertical coupling}

The tilt-calibration results left open whether gravity-referenced
communication is difficult because the sender--receiver code itself is hard to
reach, or because colonies must evolve the vertical geometry and the
sender--receiver code at the same time. To separate these possibilities, we ran
a constrained probe with the geometric precondition supplied. Initial comb tilt
was fixed near vertical at $t=0.95$, comb-tilt mutation was set to 0, comb
orientation was treated axially, the vertical-comb benefit was set to 0, and
the baseline food distribution was otherwise retained. We then varied only the
sender--receiver transposition mutation correlation and ran each condition for
120 generations over five seeds (101--105). A run was counted as reaching
gravity-referenced communication when both mean sender and receiver
transposition reached 0.50.

% Generated by experiments/run_report_artifacts.py; do not edit by hand.
% Result command: python -u experiments/run_report_artifacts.py results --only vertical_coupling_probe
% Result file: results/vertical_coupling_probe.csv
\begin{table}[h]
\centering
\small
\begin{tabular}{rrrrrrrrr}
\toprule
Corr. & Reached & Reach gen. & Prec. & Sender & Receiver & Gap & Success & Payoff \\
\midrule
0.000 & 5/5 & 42.600 & 0.878 & 0.864 & 0.881 & 0.033 & 0.192 & 0.680 \\
0.300 & 5/5 & 46.200 & 0.847 & 0.913 & 0.898 & 0.022 & 0.192 & 0.694 \\
0.600 & 5/5 & 21.000 & 0.882 & 0.890 & 0.897 & 0.010 & 0.199 & 0.754 \\
0.900 & 5/5 & 22.400 & 0.863 & 0.907 & 0.903 & 0.005 & 0.194 & 0.723 \\
1.000 & 5/5 & 23.400 & 0.836 & 0.901 & 0.901 & 0.000 & 0.191 & 0.677 \\
\bottomrule
\end{tabular}
\caption{Constrained near-vertical coupling probe. Corr. is the correlation
between sender and receiver transposition mutation increments. ``Reached''
counts seeds in which both mean sender and receiver transposition reached 0.50
by generation 120. Reach generation is averaged over reached seeds only.
Gap is the final mean absolute sender--receiver transposition difference.}
\label{tab:vertical-coupling-probe}
\end{table}


The result is strongly diagnostic. With near-vertical geometry supplied,
gravity-referenced communication is reachable in every condition. Independent
or weakly coupled sender--receiver mutations still reach the gravity channel,
but only after roughly 40--50 generations on average. Moderate to strong
coupling accelerates the transition: correlations of 0.6--1.0 reach the
threshold by about generation 21--23, while also keeping final sender and
receiver transposition closely matched. The best payoff in this probe occurs
at correlation 0.6. These results suggest that the earlier vertical-comb
failure is not because gravity-referenced communication is intrinsically
unreachable. Rather, the difficult part is the joint transition in which comb
architecture, orientation representation, and the sender--receiver code must
all become useful together.

\subsection{Long transition runs}

We therefore ran longer transition experiments using the axial orientation
representation. These runs used the baseline food distribution, 120
generations, ten seeds (101--110), sender--receiver transposition mutation
correlation 0.6, and the same three initial comb tilts and vertical-comb
benefits as the calibration grid. Unlike the constrained near-vertical probe,
comb tilt was free to mutate. A seed was counted as reaching the gravity
channel when both mean sender and receiver transposition reached 0.50. It was
counted as retaining verticality when final mean comb tilt was at least 0.80.
It was counted as collapsed when mean foraging success fell to 0.02 or below
at any generation, and as recovered when a collapsed seed ended with final
success at least 0.10.

% Generated by experiments/run_report_artifacts.py; do not edit by hand.
% Result command: python -u experiments/run_report_artifacts.py results --only long_vertical_transition
% Result file: results/long_vertical_transition_summary.csv
\begin{table}[h]
\centering
\small
\begin{tabular}{rrrrrrrr}
\toprule
Initial tilt & Benefit & Gravity & Vertical & Collapse & Recovered & Success & Final tilt \\
\midrule
0.000 & 0.150 & 0/10 & 0/10 & 0/10 & 0/10 & 0.199 & 0.118 \\
0.000 & 0.200 & 0/10 & 0/10 & 0/10 & 0/10 & 0.193 & 0.121 \\
0.000 & 0.250 & 1/10 & 1/10 & 1/10 & 0/10 & 0.175 & 0.211 \\
0.500 & 0.150 & 3/10 & 3/10 & 3/10 & 0/10 & 0.136 & 0.366 \\
0.500 & 0.200 & 4/10 & 6/10 & 6/10 & 1/10 & 0.099 & 0.599 \\
0.500 & 0.250 & 5/10 & 9/10 & 9/10 & 1/10 & 0.037 & 0.838 \\
1.000 & 0.150 & 7/10 & 9/10 & 7/10 & 1/10 & 0.079 & 0.829 \\
1.000 & 0.200 & 5/10 & 10/10 & 10/10 & 1/10 & 0.020 & 0.899 \\
1.000 & 0.250 & 8/10 & 9/10 & 9/10 & 4/10 & 0.095 & 0.829 \\
\bottomrule
\end{tabular}
\caption{Long axial-orientation transition experiment. Gravity counts seeds in
which both mean sender and receiver transposition reached 0.50 by generation
120. Vertical counts seeds with final mean comb tilt at least 0.80. Collapse
counts seeds whose mean foraging success fell to 0.02 or below at any
generation. Recovered counts collapsed seeds whose final success was at least
0.10. Success and final tilt are final-generation means over all ten seeds.}
\label{tab:long-transition}
\end{table}


The longer horizon changes the interpretation but does not by itself solve the
transition. From horizontal starts, weak and moderate vertical-comb benefits
preserve high foraging success but do not move colonies far from horizontal.
At benefit 0.25, one seed reaches verticality and the gravity channel, but that
same seed collapses. From tilted starts, higher benefits increasingly retain
verticality and allow more seeds to reach the gravity channel, but collapse
also becomes common; at benefit 0.20, one seed recovers to a high-success
vertical gravity-channel state. From vertical starts, longer runs show that
successful vertical gravity-channel outcomes are possible, especially at
benefit 0.25 where four collapsed seeds recover by generation 120, but most
seeds still experience collapse. Thus the 30-generation failures were too short
to distinguish an impossible transition from a slow coordination process, but
120 generations show that longer time alone is insufficient to make the
transition robust under the current parameters.

\section{Conclusion}

The current model now separates world directions, comb-plane dance signals,
comb tilt and orientation, and the sun/gravity reference. Direct pointing is
available on horizontal combs without additional mapping machinery, while the
sun/gravity channel becomes useful only as gravity projects into the comb
plane. This supports the intended conceptual setup: a non-communication benefit
for vertical combs can create pressure toward verticality, after which
sender--receiver transposition may become advantageous.

The calibration runs show that this transition is not yet smooth. When
verticality is weakly rewarded, colonies starting from horizontal combs remain
mostly horizontal. When verticality is strongly rewarded, vertical combs can be
retained, but foraging success often collapses before communication adapts.
Orientation is not yet robustly selected relative to the sun arc under the
circular representation. The constrained vertical-coupling probe clarifies this
picture: once near-vertical geometry is supplied and orientation is treated
axially, gravity-referenced sender--receiver transposition evolves reliably,
especially when sender and receiver transposition mutations are moderately or
strongly coupled. The long axial-orientation transition runs show that the
30-generation failures were partly a horizon problem: some tilted or vertical
starts do recover into high-success vertical gravity-channel states by
generation 120. However, the transition is still not robust. Horizontal starts
mostly remain horizontal unless the vertical-comb benefit is strong enough to
risk collapse, and vertical-retaining seeds often lose foraging performance
before communication catches up. The next modeling priority is therefore to
calibrate vertical-comb benefits, mutation scales, and coupling assumptions so
that architecture and communication can coevolve without routinely passing
through a collapse in foraging success.

\bibliographystyle{plain}
\bibliography{references}

\end{document}
